{\footnotesize
    \notation $\displaystyle\bar{X}_n = \frac{1}{n}S_n = \frac{1}{n}\sum_{i=1}^{n}X_n$
}

\theorem \textbf{Schwaches Gesetz der grossen Zahlen}\\
\smalltext{$X_1,X_2,\ldots$ unabh.,$\quad\forall k: \E[X_k] = \mu,\V[X_k]=\sigma^2$}

Es konvergiert $\bar{X}_n$ für $n\to\infty$ gegen $\mu = \E[X_i]$ 
$$
    \P\Bigl[ |\bar{X}_n - \mu| < \epsilon \Bigr] \overset{n\to\infty}{\rightarrow} 0
$$
{\footnotesize
    \remark $\V\Bigl[\overbrace{\frac{1}{n}\sum_{i=1}^nX_i}^{\bar{X}_n}\Bigr] = \frac{1}{n^2}\displaystyle\sum_{i=1}^n\V[X_i] = \frac{n}{n^2}\V[X_i] = \frac{1}{n}\V[X_i]$
}

\theorem \textbf{Starkes Gesetz der grossen Zahlen}\\
\smalltext{$X_1,X_2,\ldots$ unabh.,$\quad\forall k: \E[X_k] = \mu$}
$$
    \P\biggl[ \Bigl\{ \omega\in\Omega \sep \bar{X}_n \overset{n\to\infty}{\to}\mu \Bigr\} \biggr] = 1
$$

\definition \textbf{Konvergenz in Verteilung}\\
\smalltext{$(X_n)_{n\in\mathbb{N}}, X$ Zufallsvariablen}
$$
    X_n \overset{\text{Approx.}}{\approx} X \qquad \text{für } n \to \infty
$$
Falls:
$$
    \forall x \in \R:\quad \underset{n\to\infty}{\lim}\P\Bigl[ X_n \leq x \Bigr] = \P\Bigl[X \leq x\Bigr]
$$

{\footnotesize
    \remark Diskrete Zufallsvariablen können zu stetigen konvergieren.
}




% Slides viel detaillierter als Skript
\theorem \textbf{Zentraler Grentwertsatz}\\
\smalltext{$(X_k)_{k\geq1}$ i.i.d s.d. $\E[X_k]=\mu, \V[X_k]=\sigma^2, S_n = \sum_{i=1}^{n}X_i$}
$$
    \underset{n\to\infty}{\lim}\P\Biggl[ \frac{S_n - n\mu}{\sigma\sqrt{n}} \leq x \Biggr] = \Phi(x)
$$
$\Phi(x)$ ist die Verteilung von $\mathcal{N}(0, 1)$.
$$
    S_n \overset{\text{approx.}}{\sim} \mathcal{N}(n\mu, n\sigma^2) 
$$
\subtext{ZGS wird häufig zur Approximation von Summen verwendet}

\definition \textbf{Standardisierung von $S_n$}
$$
    S_n^* := \frac{S_n - n\mu}{\sigma\sqrt{n}} = \frac{S_n - \E[S_n]}{\sqrt{\V[S_n]}} \sim \mathcal{N}(0,1)
$$
\subtext{Im Skript auch $Z_n$}

% Slides: Chernoff bounds & mehr ZGS (nicht im Skript)

\newpage

{\footnotesize
    \textbf{Beispiel}: Approximation via Standardisierung mit ZGS.

    Sei $X_i,\ldots,X_n \overset{\text{i.i.d.}}{\sim} \P_X$ s.d. $\forall i: \E[X_i]=8, \V[X_i]=1.44$.

    für $S_{100} = \sum_{i=1}^{n}X_i$ approximieren wir $\P\Bigl[ 788 \leq S_100 \leq 824 \Bigr]$\\
    Es gilt: $\E[S_{100}] = 100\cdot\E[X_i] = 800,\quad \V[S_{100}]=100\cdot\V[X_i] = 144$
    \begin{align*}
            &\P\Bigl[ 788 \leq S_100 \leq 824 \Bigr] \\
        =   &\P\Biggl[\frac{788 - \E[S_{100}]}{\sqrt{\V[S_{100}]}} \leq \frac{S_{100} - \E[S_{100}]}{\sqrt{\V[S_{100}]}} \leq \frac{824 - \E[S_{100}]}{\sqrt{\V[S_{100}]}}\Biggr] \\
        =   &\P\Bigl[ -1 \leq \frac{S_{100} - \E[S_{100}]}{\sqrt{\V[S_{100}]}} \leq 2 \Bigr] \\
        =   & \Phi(2) - \Phi(-1)    \\
        =   & \Phi(2) - (1-\Phi(1)) \\
        =   & \Phi(2) + \Phi(1) - 1 \\
        \approx & 0.8185   
    \end{align*} 
    Genutztz wurde die Symmetrie von $\Phi$ und $\frac{S_{100} - \E[S_{100}]}{\sqrt{\V[S_{100}]}} \sim \mathcal{N}(0,1)$.\\
    \subtext{Gemäss ZGS Nur eine approximation, da wir hier $n=100$ nutzen.} 
}

{\footnotesize
    \textbf{Beispiel}: Zentraler Grenzwertsatz als Teststatistik

    $X_1,\ldots,X_{72} \overset{\text{i.i.d.}}{\sim} \text{Ber}(\theta)$, $S_{72} = \sum_{i=1}^{72} X_i$. Stichprobe $\bar{s}_{72}=32$.
    $$
        H_0: \theta = \frac{1}{3} \qquad H_A: \theta > \frac{1}{3}
    $$
    Mit ZGS, wobei wir $\theta = \frac{1}{3}$ ($H_0$) nutzen, erhalten wir $T$:
    $$
        T \overset{\text{ZGS}}{=} \frac{S_n - n\mu}{\sigma\sqrt{n}} = \frac{S_n - n\E[X_i]}{\sqrt{\V[X_i]n}} \overset{\text{Ber.}}{=} \frac{S_n - 72\cdot\frac{1}{3}}{\sqrt{\frac{2}{9}\cdot 72}} = \boxed{\frac{S_n-24}{4}}
    $$
    Der Vorteil hiervon ist, der ZGS garantiert: $T \sim \mathcal{N}(n\mu, n\sigma^2)$:

    Realisierungen $T(\omega) = T(\bar{s}_n) = T(32) = 2$ bilden mit $z_{1-\alpha}$: $K = (z_{1-\alpha},\infty)$ und können direkt ausgewertet werden, z.B. für $p$-Wert.

    Hier z.B. $p \approx \P[Z \geq 2] = 1 - \Phi(2)\approx 0.023$, $Z \sim \mathcal{N}(0,1)$.\\
    $p=0.023>0.01$ bedeutet, wir akzeptieren $H_0$ auf dem $\alpha=0.01$ (1\%) Niveau.
}